\chapter*{Úvod}
\addcontentsline{toc}{chapter}{Úvod}
\thispagestyle{fancy}
\fancyfoot[R]{\thepage}

{\color{infoBlue}
(obsahuje cíl práce, motivace k výběru tématu, aktuálnost apod.)

Závěrečná práce je psaná písmem Calibri (Základní text) velikost 11 b., řádkování je 1,15 a mezera za odstavci je 6 b. Pro zvýrazňování v textu je nejvhodnější \textit{kurzíva}, omezeně lze použít \textbf{tučné písmo}. Nikdy se nepoužívá \underline{podtrhávání}.

Každá kapitola (nadpis 1. úrovně) začíná na novém listě.

U závěrečné práce se předpokládá, že tištěná jednostranně a vázaná, proto je velikost okrajů nastavena na 2,5 cm nahoře, dole a vpravo; 3,5 cm vlevo. Práce bude vytištěna jednostranně, barevně. 

Tato šablona slouží jako šablona pro psaní závěrečné práce na VŠPJ. Nejedná se o návod, jak práci zpracovávat. V dokumentu Jak psát práce na VŠPJ: Typografická pravidla pro studenty VŠPJ (Borůvková a kol., 2021) lze najít přehled nejdůležitějších pravidel pro formátování závěrečné práce na VŠPJ včetně typografických pravidel, ve kterých se nejčastěji chybuje, a ukázek citace použité literatury.

Skladba a názvy kapitol jsou doporučené. Podle zpracovávaného tématu a zvyklostí odborných kateder mohou být upraveny.
}

Vlastní text úvodu. Vlastní text úvodu. Vlastní text úvodu. Vlastní text úvodu. Vlastní text úvodu. Vlastní text úvodu. Vlastní text úvodu. Vlastní text úvodu. Vlastní text úvodu. Vlastní text úvodu. Vlastní text úvodu.

\pagebreak
